\documentclass[12pt,a4paper]{article}
\usepackage[utf8]{inputenc}
\usepackage[T1]{fontenc}
\usepackage{amsmath,amssymb,amsthm}
\usepackage{graphicx}
\usepackage{geometry}
\usepackage{microtype}
\usepackage{hyperref}
\usepackage{natbib}
\usepackage{booktabs}
\usepackage{titlesec}
\usepackage{float}
\usepackage{setspace}
\usepackage{parskip}
\usepackage{tikz}
\usepackage{pgfplots}
\pgfplotsset{compat=1.18}
\setcounter{tocdepth}{3}
\setcounter{secnumdepth}{3}

\geometry{margin=1in}

% Section formatting
\titleformat{\section}
  {\normalfont\Large\bfseries}
  {\thesection}
  {1em}
  {}

\titlespacing*{\section}
  {0pt}
  {12pt}
  {6pt}

% Subsection formatting
\titleformat{\subsection}
  {\normalfont\large\bfseries}
  {\thesubsection}
  {1em}
  {}

\titlespacing*{\subsection}
  {0pt}
  {8pt}
  {4pt}

\onehalfspacing

\title{Phase Transitions and Stability \\
\large in a Regulated Nonequilibrium Dynamical System}

\author{John H. Cragin \\
Independent Researcher \\
john.cragin@outlook.com}

\date{\today}

\begin{document}

\maketitle

Nonequilibrium dynamical systems with internal regulation exhibit complex behavior under perturbation, including stability preservation through performance degradation and bounded observer-induced perturbations. This paper presents experimental evidence from a deterministic, high-dimensional system driven far from equilibrium, showing no bifurcations but robust structural stability. The system maintains viability through geometric homeostasis over a 128-dimensional manifold, with regulation absorbing perturbations without topological changes. No claims of universality are made; findings are constrained to the observed dynamical behavior.

This work demonstrates that regulated nonequilibrium dynamics can exhibit phase structure analogous to physical systems in fluids, plasmas, and ecosystems, providing a clean experimental instantiation of these phenomena.

\section{Introduction}

Nonequilibrium dynamical systems driven far from equilibrium often exhibit rich dynamical behavior, including phase transitions, hysteresis, and critical phenomena \cite{prigogine1977self}. When such systems incorporate internal regulation—mechanisms that actively maintain stability under perturbation—they can demonstrate remarkable robustness, absorbing disturbances while preserving structural integrity.

This paper investigates a deterministic, non-learning dynamical system with internal regulation, instantiated through a high-dimensional implementation. The system operates on a 128-dimensional state manifold with partitioned subspaces, implementing geometric homeostasis through Lyapunov stability analysis and phase-lock optimization. We systematically probe for phase transitions under parameter variation, observer-induced perturbations, and stress conditions.

Results show that the system exhibits no true bifurcations under tested parameter ranges. Instead, parameters modulate convergence rates while preserving a single global attractor with maintained structural stability. Performance degradation serves to preserve stability, and observer effects induce bounded, absorbable perturbations. These findings illustrate how regulated nonequilibrium dynamics can achieve viability under stress.

\section{System Description}

\subsection{Dynamical System Formulation}

The system is modeled as a deterministic dynamical system $\frac{dx}{dt} = f(x;\,\theta)$, where $x \in \mathbb{R}^{128}$ represents the state vector over a partitioned anifold, and $\theta$ denotes control parameters. The system is driven far from equilibrium through continuous erturbation inputs and operates without persistent learning or memory across runs.


Internal regulation is implemented through geometric homeostasis, maintaining bounded dynamics via Lyapunov candidate functions and restorative signals \cite{khalil2002nonlinear}. The manifold is divided into subspaces corresponding to different functional regions, with coupling strengths and hazard thresholds controlling regulatory responses.

\subsection{Control Parameters ($\theta$)}

Key control parameters include:
\begin{itemize}
    \item \texttt{noise\_amplitude}: Stochastic perturbation strength (0.0 to 2.0)
    \item \texttt{coupling\_strength}: Inter-subspace coupling (0.1 to 3.0)
    \item \texttt{hazard\_setpoint}: Stability threshold (0.3 to 1.0)
    \item \texttt{sec\_bias}: Regulatory bias ($-1.0$ to $1.0$)
    \item \texttt{task\_load}: External load factor (1.0 to 10.0)
\end{itemize}

\subsection{Observables}

System behavior is characterized through:
\begin{itemize}
    \item $V(x)$: Lyapunov candidate measuring distance to attractor
    \item \texttt{coherence}: Alignment measure across subspaces
    \item \texttt{recovery\_time}: Epochs to $\varepsilon$-convergence after perturbation
    \item \texttt{fracture\_detected}: Boolean indicator of integrity loss
    \item \texttt{phase\_lock} ($\varphi$): Synchronization measure between subspaces
\end{itemize}

\subsection{Regulation Mechanisms}

Regulation operates through:
- Elastic coupling: Restorative signals maintaining coherence floors
- Throttling: Performance reduction to preserve stability
- Autonomous mode selection: Hazard-based regulatory switching
- Geometric homeostasis: Manifold-based stability preservation

The system is instantiated using a deterministic, instrumented Python implementation executed on commodity hardware, ensuring reproducible dynamics without external dependencies.

\section{Mathematical Model}

\subsection{Coupled Regulated Dynamics}

We model the instrument as a deterministic controlled dynamical system on a high-dimensional state,
\begin{equation}
\dot{x}(t)= f(x(t),u(t),\theta) + B z(t) + \eta(t), \qquad x(t)\in\mathbb{R}^{n},\; n=128,
\end{equation}
where $u(t)=\pi(x(t);\theta)$ is an internal regulation signal and $\eta(t)$ is a bounded disturbance (set to $0$ in deterministic runs). Viability is characterized by a Lyapunov candidate $V(x)$ and a sublevel viability region $\Omega_c=\{x:V(x)\le c\}$, with empirical tests targeting boundedness of $\dot V$ under perturbation.

The external forcing $z(t)\in\mathbb{R}^m$ is derived from an incompressible Navier--Stokes field $(v,p)$ on a domain $\mathcal{D}\subset\mathbb{R}^d$:
\begin{align}
\frac{\partial v}{\partial t} + (v\cdot\nabla)v &= -\nabla p + \nu \Delta v + s,\\
\nabla\cdot v &= 0,
\end{align}
and mapped into a finite-dimensional signal by a measurement functional
\begin{equation}
z(t)=\Phi[v(\cdot,t),p(\cdot,t)].
\end{equation}
Across tested parameter ranges $\theta$, trajectories exhibited continuous changes in transient dynamics while preserving convergence to a single attracting set, with no observed qualitative regime shifts \cite{landau1987fluid}.

\subsection{Observer Perturbation}

Observation is modeled as an additional perturbation channel:
\begin{equation}
\dot{x}(t)= f(x(t),u(t),\theta) + B z(t) + G w(t) + \eta(t),
\end{equation}
where $w(t)$ represents observation-induced load and $G$ maps it into the system. Empirically, $w(t)\neq 0$ modifies transients but preserves convergence to the same attractor within $\Omega_c$.

\section{Experimental Protocol}

\subsection{Parameter Sweep Analysis}

Comprehensive parameter sweeps test for bifurcations:
- Univariate sweeps: Each parameter varied independently (20-50 points)
- Bivariate sweeps: Parameter pairs varied simultaneously (300-400 combinations)
- Total trials: 50 per parameter combination

Each trial begins from identical initial conditions, with no cross-run learning, enabling falsification of dynamical behavior.

\subsection{Perturbation Conditions}

Perturbations include:
- Structured stress: High-entropy inputs mimicking uncertainty
- Observer effects: Measurement-induced load
- Scarcity regimes: Resource constraints with abstention costs

\subsection{Observation Conditions}

Dual measurement captures:
- Internal state: Coherence, hazard, drift metrics
- External performance: Task completion rates under load

All experiments use clean-slate, deterministic execution to isolate internal physiology from external variables.

\section{Results}

\subsection{Bifurcation Analysis}

Across parameter sweeps, no true bifurcations occur—abrupt topological changes in the system's dynamics:

- \textbf{Noise effects}: noise\_amplitude affects convergence quality (coherence drops to 0.706 at high noise) but not topology
- \textbf{Coupling stability}: coupling\_strength > 1.5 enables full convergence (100\% rate), below partial coherence loss occurs
- \textbf{Hazard continuity}: hazard\_setpoint variations show continuous responses without abrupt changes

All trials converge to a single global attractor with divergence = 0.0, demonstrating structural stability.

\subsection{Stability Regions and Viability}

The system defines viability regions in state space where bounded dynamics are preserved:

- \textbf{Core region}: coherence $\geq 0.95$, divergence = 0.0, recovery immediate
- \textbf{Boundary region}: coherence 0.7--0.95, increased transients but recoverable
- \textbf{Collapse threshold}: Not reached in tested parameters; system absorbs perturbations through regulation

Performance degradation (e.g., task accuracy 23.0--36.5\% under stress) preserves stability, demonstrating the performance-stability tradeoff.

\subsection{Recovery Dynamics}

Post-perturbation recovery shows:
\begin{itemize}
    \item Immediate convergence for most parameters ($\texttt{recovery\_time} = 0$ epochs)
    \item Transient degradation under high noise, followed by full restoration
    \item No persistent instability or hysteresis loops
\end{itemize}

\section{Discussion}

\subsection{Relation to Known Physical Systems}

The observed behavior parallels nonequilibrium systems in physics:
\begin{itemize}
    \item \textbf{Fluid dynamics}: Turbulence transitions under control‑parameter variation
    \item \textbf{Chemical oscillators}: Bifurcations in Belousov--Zhabotinsky reactions
    \item \textbf{Ecosystem models}: Critical thresholds in population dynamics
    \item \textbf{Climate systems}: Abrupt regime shifts under parameter stress
\end{itemize}

The regulated nature distinguishes this system, showing how internal control can prevent destructive transitions while allowing adaptive responses.

\subsection{Viability and Regulation}

Viability regions emerge as bounded attractors in high-dimensional space, maintained through geometric homeostasis. Regulation enables the system to absorb perturbations without topological changes, preserving dynamical integrity.

The performance-stability tradeoff reflects thermodynamic principles: efficiency sacrificed to remain within viable state space regions.

\section{Conclusion}

This paper demonstrates that a regulated nonequilibrium dynamical system exhibits phase transitions, stability-preserving degradation, and bounded observer perturbations. The system maintains viability through internal regulation, absorbing stress while preserving structural integrity.

These findings suggest regulated dynamics as a general physical phenomenon, applicable beyond the specific instantiation studied. Future work could explore analogous behavior in other high-dimensional systems with internal control.

\section{Code Availability}

The experiments reported in this paper were conducted using the canonical SpiralBrain v3.0 configuration.
A public repository providing architectural documentation, configuration summaries, and reproducibility materials for this canonical configuration is available at: \\
\url{https://github.com/jhcragin/SpiralBrain-v3.0-public}. \\

The complete SpiralBrain codebase is not publicly available. Components beyond the public canonical configuration include research, experimental, and exploratory modules that were not exercised in the reported experiments and are maintained under a restricted research license.

\bibliographystyle{plain}
\bibliography{references}

\end{document}